% ============================================================
\chapter{Mod\`eles Bay\'esiens Non-Param\'etriques}
\label{ch:non_parametrique}
% ============================================================

En statistique param\'etrique, on choisit \`a l'avance le nombre de param\`etres~: un mod\`ele de m\'elange \`a $K$ composantes, une r\'egression polynomiale de degr\'e $p$. Mais comment choisir $K$ ou $p$~? La statistique bay\'esienne non param\'etrique contourne le probl\`eme en pla\c{c}ant des lois a priori sur des espaces de \emph{dimension infinie}~: le processus de Dirichlet (Ferguson, 1973) d\'efinit une loi sur l'espace des distributions de probabilit\'e~; le processus gaussien d\'efinit une loi sur l'espace des fonctions. Le mod\`ele adapte alors automatiquement sa complexit\'e aux donn\'ees~: le nombre de clusters cro\^it avec la taille de l'\'echantillon, la souplesse de la r\'egression s'ajuste. Cette flexibilit\'e a r\'evolutionn\'e le clustering, la mod\'elisation de donn\'ees de survie et l'analyse de textes.

\begin{intuition}[title=\textbf{Id\'ee centrale}]
En statistique bay\'esienne non param\'etrique, on place des lois a priori
sur des espaces de dimension infinie (distributions, fonctions, partitions).
Le mod\`ele adapte sa complexit\'e aux donn\'ees~: le nombre de composantes
ou la souplesse de la fonction ajust\'ee cro\^it avec $n$.
\end{intuition}

% ------------------------------------------------------------
\section{Motivation}
% ------------------------------------------------------------

Les mod\`eles param\'etriques fixent le nombre de param\`etres \`a l'avance~:
un m\'elange de $K$ gaussiennes, une r\'egression d'ordre $p$, etc.
Or, le choix de $K$ ou $p$ est souvent arbitraire.

\begin{remarque}
Le terme \og non param\'etrique \fg est trompeur~: il y a bien des
param\`etres, mais en nombre \emph{infini} (ou croissant avec les donn\'ees).
\end{remarque}

% ------------------------------------------------------------
\section{Le processus de Dirichlet}
% ------------------------------------------------------------

\begin{definition}[Processus de Dirichlet]
\index{processus de Dirichlet}
Soit $\alpha > 0$ un param\`etre de concentration et $G_0$ une mesure de
probabilit\'e de base sur $(\Theta, \mathcal{B})$. On dit que
$G \sim \mathrm{DP}(\alpha, G_0)$ si, pour toute partition mesurable
$(A_1, \ldots, A_K)$ de $\Theta$~:
\[
  \bigl(G(A_1), \ldots, G(A_K)\bigr) \sim
  \mathrm{Dir}\!\bigl(\alpha G_0(A_1), \ldots, \alpha G_0(A_K)\bigr).
\]
\end{definition}

\begin{proposition}[Propri\'et\'es du DP]
Soit $G \sim \mathrm{DP}(\alpha, G_0)$. Alors~:
\begin{enumerate}
  \item $\E[G(A)] = G_0(A)$ pour tout $A$.
  \item $\Var(G(A)) = \dfrac{G_0(A)(1 - G_0(A))}{\alpha + 1}$.
  \item $G$ est presque s\^urement discr\`ete (mesure atomique).
\end{enumerate}
\end{proposition}

\begin{attention}
Le DP g\'en\`ere des distributions discr\`etes m\^eme si $G_0$ est continue.
C'est cette discr\'etude qui induit le regroupement (\emph{clustering}).
\end{attention}

% ------------------------------------------------------------
\section{Construction par brisure de b\^aton}
% ------------------------------------------------------------

\begin{theoreme}[Sethuraman, 1994]
\index{brisure de b\^aton}
Si $G \sim \mathrm{DP}(\alpha, G_0)$, alors $G$ admet la repr\'esentation~:
\[
  G = \sum_{k=1}^{\infty} \pi_k \,\delta_{\theta_k^*},
\]
o\`u $\theta_k^* \stackrel{\text{iid}}{\sim} G_0$ et les poids sont
construits par brisure de b\^aton~:
\[
  V_k \stackrel{\text{iid}}{\sim} \mathrm{Beta}(1, \alpha), \quad
  \pi_k = V_k \prod_{j=1}^{k-1}(1 - V_j).
\]
\end{theoreme}

\begin{remarque}
Les poids $\pi_k$ d\'ecroissent stochastiquement~: les premiers atomes
re\c{c}oivent l'essentiel de la masse. En pratique, on tronque \`a $K$
termes pour l'inf\'erence.
\end{remarque}

\begin{minted}{python}
import numpy as np

def stick_breaking(alpha, K):
    """Genere K poids par brisure de baton."""
    V = np.random.beta(1, alpha, size=K)
    pi = np.zeros(K)
    remaining = 1.0
    for k in range(K):
        pi[k] = V[k] * remaining
        remaining *= (1 - V[k])
    return pi
\end{minted}

% ------------------------------------------------------------
\section{Le processus du restaurant chinois}
% ------------------------------------------------------------

\begin{definition}[Processus du restaurant chinois (CRP)]
\index{processus du restaurant chinois}
Soit $(\theta_1, \theta_2, \ldots)$ un \'echantillon du DP marginal.
La loi conditionnelle pr\'edictive est~:
\[
  \theta_{n+1} \mid \theta_1, \ldots, \theta_n \sim
  \frac{1}{n + \alpha}\left(\sum_{k=1}^{K_n} n_k\,\delta_{\theta_k^*}
  + \alpha\,G_0\right),
\]
o\`u $\theta_1^*, \ldots, \theta_{K_n}^*$ sont les valeurs distinctes et $n_k$
leurs effectifs.
\end{definition}

\begin{intuition}[title=\textbf{M\'etaphore du restaurant}]
Des clients arrivent un par un dans un restaurant avec des tables infinies.
Le client $n+1$ s'assoit \`a la table $k$ (d\'ej\`a occup\'ee par $n_k$ clients)
avec probabilit\'e $n_k / (n + \alpha)$, ou ouvre une nouvelle table avec
probabilit\'e $\alpha / (n + \alpha)$.
\end{intuition}

\begin{proposition}[Nombre attendu de clusters]
Pour $n$ observations tir\'ees d'un $\mathrm{DP}(\alpha, G_0)$, le nombre
attendu de clusters distincts est~:
\[
  \E[K_n] = \alpha \sum_{i=1}^{n} \frac{1}{\alpha + i - 1}
           \approx \alpha \ln\!\left(1 + \frac{n}{\alpha}\right).
\]
\end{proposition}

% ------------------------------------------------------------
\section{M\'elange par processus de Dirichlet (DP-GMM)}
% ------------------------------------------------------------

\begin{definition}[Mod\`ele DP-GMM]
\index{DP-GMM}
Le \textbf{m\'elange gaussien infini} (DP-GMM) est d\'efini par~:
\begin{align}
  G &\sim \mathrm{DP}(\alpha, G_0), \\
  (\mu_i, \Sigma_i) &\sim G, \quad i = 1, \ldots, n, \\
  x_i \mid \mu_i, \Sigma_i &\sim \mathcal{N}(\mu_i, \Sigma_i).
\end{align}
Le nombre de composantes est d\'etermin\'e par les donn\'ees.
\end{definition}

\begin{exemple}
Avec $G_0 = \mathrm{NIW}(\mu_0, \kappa_0, \nu_0, \Psi_0)$ (Normal-Inverse-Wishart),
l'\'echantillonneur de Gibbs du CRP it\`ere~:
\begin{enumerate}
  \item Pour chaque $i$, r\'eassigner $z_i$ selon $p(z_i = k \mid \ldots) \propto
        n_{-i,k}\,\mathcal{N}(x_i \mid \mu_k, \Sigma_k)$ ou cr\'eer un nouveau cluster.
  \item Pour chaque cluster $k$, mettre \`a jour $(\mu_k, \Sigma_k)$ via le posterior NIW.
\end{enumerate}
\end{exemple}

\begin{minted}{python}
from sklearn.mixture import BayesianGaussianMixture

# DP-GMM via scikit-learn (approximation variationnelle)
dpgmm = BayesianGaussianMixture(
    n_components=20,            # borne superieure
    covariance_type='full',
    weight_concentration_prior_type='dirichlet_process',
    weight_concentration_prior=0.1
)
dpgmm.fit(X)
labels = dpgmm.predict(X)
\end{minted}

% ------------------------------------------------------------
\section{Processus gaussiens comme priors sur les fonctions}
% ------------------------------------------------------------

\begin{definition}[Processus gaussien]
\index{processus gaussien}
Un \textbf{processus gaussien} (GP) est une collection de variables
al\'eatoires $\{f(x) : x \in \mathcal{X}\}$ telle que toute sous-collection
finie suit une loi gaussienne multivari\'ee~:
\[
  f \sim \mathcal{GP}(m, k), \quad
  (f(x_1), \ldots, f(x_n)) \sim \mathcal{N}\!\bigl(\mathbf{m}, \mathbf{K}\bigr),
\]
o\`u $m_{i} = m(x_i)$ et $K_{ij} = k(x_i, x_j)$.
\end{definition}

\begin{theoreme}[Posterior du GP pour la r\'egression]
Soit $\mathbf{y} = f(\mathbf{X}) + \boldsymbol{\epsilon}$,
$\boldsymbol{\epsilon} \sim \mathcal{N}(\mathbf{0}, \sigma^2 \mathbf{I})$.
La loi a posteriori de $f_* = f(x_*)$ est~:
\begin{align}
  f_* \mid \mathbf{X}, \mathbf{y}, x_* &\sim \mathcal{N}(\bar{f}_*, \Var(f_*)), \\
  \bar{f}_* &= \mathbf{k}_*^\top (\mathbf{K} + \sigma^2 \mathbf{I})^{-1}\mathbf{y}, \\
  \Var(f_*) &= k(x_*, x_*) - \mathbf{k}_*^\top(\mathbf{K} + \sigma^2 \mathbf{I})^{-1}\mathbf{k}_*.
\end{align}
\end{theoreme}

\begin{remarque}
Le choix du noyau $k$ encode les hypoth\`eses sur la r\'egularit\'e de $f$~:
RBF pour des fonctions lisses, Mat\'ern pour un contr\^ole fin de la
diff\'erentiabilit\'e, p\'eriodique pour des donn\'ees cycliques.
\end{remarque}

% ------------------------------------------------------------
\section{Autres processus non param\'etriques}
% ------------------------------------------------------------

\begin{definition}[Processus du buffet indien (IBP)]
\index{processus du buffet indien}
Le \textbf{IBP} est un prior sur des matrices binaires $Z$ de taille
$n \times \infty$ utilis\'e pour la s\'election de caract\'eristiques.
Chaque client (observation) choisit des plats (caract\'eristiques) existants
avec probabilit\'e $m_k / n$ et essaie $\mathrm{Poisson}(\alpha / n)$
nouveaux plats.
\end{definition}

\begin{definition}[Processus de Pitman--Yor]
\index{processus de Pitman--Yor}
Le \textbf{processus de Pitman--Yor} $\mathrm{PY}(d, \alpha, G_0)$ g\'en\'eralise
le DP avec un param\`etre de discount $d \in [0,1)$. La loi pr\'edictive est~:
\[
  \theta_{n+1} \mid \theta_{1:n} \sim
  \frac{1}{n + \alpha}\left(\sum_{k=1}^{K_n}(n_k - d)\,\delta_{\theta_k^*}
  + (\alpha + d K_n)\,G_0\right).
\]
Pour $d = 0$, on retrouve le DP.
\end{definition}

% ------------------------------------------------------------
\section{Formules cl\'es}
% ------------------------------------------------------------

\begin{formules}
\begin{align}
  \text{DP :} \quad & (G(A_1), \ldots, G(A_K)) \sim \mathrm{Dir}(\alpha G_0(A_1), \ldots, \alpha G_0(A_K)) \\
  \text{CRP :} \quad & p(\theta_{n+1} = \theta_k^*) = \frac{n_k}{n + \alpha}, \quad
                        p(\text{nouveau}) = \frac{\alpha}{n + \alpha} \\
  \text{Brisure :} \quad & \pi_k = V_k \prod_{j<k}(1-V_j), \quad V_k \sim \mathrm{Beta}(1,\alpha) \\
  \text{GP :} \quad & \bar{f}_* = \mathbf{k}_*^\top(\mathbf{K}+\sigma^2\mathbf{I})^{-1}\mathbf{y}
\end{align}
\end{formules}

% ------------------------------------------------------------
\section{Exercices}
% ------------------------------------------------------------

\begin{exercice}
V\'erifier que les poids de la brisure de b\^aton somment \`a 1 presque
s\^urement. \emph{Indication~:} calculer $\E\bigl[\sum_{k=1}^K \pi_k\bigr]$ et
passer \`a la limite.
\end{exercice}

\begin{exercice}
Simuler le processus du restaurant chinois pour $\alpha \in \{0.1, 1, 10, 100\}$
et $n = 500$. Tracer le nombre de clusters $K_n$ en fonction de $\alpha$.
Comparer avec la formule th\'eorique $\E[K_n] \approx \alpha \ln(1 + n/\alpha)$.
\end{exercice}

\begin{exercice}
Impl\'ementer un \'echantillonneur de Gibbs pour le DP-GMM sur des donn\'ees
simul\'ees en dimension 2 (trois clusters gaussiens). Visualiser les
affectations au fil des it\'erations.
\end{exercice}

\begin{exercice}
Consid\'erer un GP avec noyau RBF $k(x,x') = \sigma_f^2 \exp(-\norm{x-x'}^2/(2\ell^2))$.
Montrer que la loi pr\'edictive du GP converge vers un interpolateur exact
lorsque $\sigma^2 \to 0$. Illustrer graphiquement avec des donn\'ees simul\'ees.
\end{exercice}

\begin{exercice}
Montrer que le processus de Pitman--Yor engendre un nombre de clusters
$K_n = O(n^d)$, contre $K_n = O(\ln n)$ pour le DP. Quelle implication
pour la mod\'elisation de donn\'ees textuelles suivant la loi de Zipf~?
\end{exercice}
